\section{Tối ưu cấu trúc dữ liệu và thuật toán}

Thách thức lớn nhất đến từ cơ chế flood fill. Khi một ô trống (giá trị 0) được mở, hệ thống phải tiếp tục mở toàn bộ các ô trống liên thông với nó. Nếu xử lý không tối ưu, thao tác này có thể gây chậm hoặc thậm chí lỗi chương trình.

\textbf{Vấn đề giới hạn đệ quy:}
Nếu dùng DFS đệ quy trong Python, nguy cơ gặp lỗi RecursionError rất có thể xảy ra, đặc biệt với bàn chơi 16×16. Một chuỗi ô trống dài liên tiếp có thể khiến số lần gọi hàm vượt quá giới hạn cho phép.

Giải pháp là chuyển sang DFS dạng lặp, sử dụng một Stack để tự quản lý quá trình duyệt. Cách này giúp kiểm soát bộ nhớ tốt hơn và tránh hoàn toàn phụ thuộc vào call stack của ngôn ngữ.

\textbf{Giảm độ trễ tính toán:}
Nếu mỗi lần duyệt một ô mới lại phải tính lại số mìn xung quanh, chương trình sẽ tốn thêm thời gian không cần thiết. Vì vậy, ngay sau khi sinh mìn, hệ thống sẽ tính trước toàn bộ số ô lân cận và lưu lại vào một ma trận riêng.

Nhờ đó, khi DFS chạy, việc truy xuất thông tin chỉ mất thời gian hằng số O(1), giúp cơ chế mở lan hoạt động mượt và ổn định.